\documentclass[../../main.tex]{subfiles}
\graphicspath{{../../image/}{../../figure/}}
\begin{document}

\section{极后向特征线}\label{section:ch10-极后向特征线}

参照某点 $(\bar{x},\bar{t})\in(-\infty,\infty)\times[0,\infty)$，过 $(\bar{x},\bar{t})$ 的广义特征线若定义在 $[0,\bar{t}]$ 上，则称为\emph{后向的}；若定义在 $[\bar{t},\infty)$ 上，则称为\emph{前向的}。过 $(\bar{x},\bar{t})$ 的极小、极大后向与前向广义特征线均以极端速度传播，因而自然地成为接触线的候选者。由 Lax $E$ 条件可知，这至少对极后向特征线成立：

\begin{theorem}\label{theorem:ch10-极后向特征线为接触}
与 $U$ 相联系、由上半平面任一点 $(\bar{x},\bar{t})$ 发出的极小（相应地极大）后向 $i$-特征线，是一条左（相应地右）$i$-接触。
\end{theorem}

\begin{proof}
设 $\xi(\cdot)$ 表示由 $(\bar{x},\bar{t})$ 发出、定义在 $[0,\bar{t}]$ 上的极小后向 $i$-特征线。固定 $\varepsilon>0$，并按下述算法选取 $\bar{t}=\tau_{0}>\tau_{1}>\cdots>\tau_{k}=0$，其中 $k\geqslant1$：先取 $\tau_{0}=\bar{t}$。设 $\tau_{m}>0$ 已确定，令 $\xi_{m}(\cdot)$ 表示由点 $(\xi(\tau_{m})-\varepsilon,\tau_{m})$ 发出的极小后向 $i$-特征线。若对 $0<t\leqslant\tau_{m}$ 有 $\xi_{m}(t)<\xi(t)$，则置 $\tau_{m+1}=0$、$m+1=k$ 并终止；否则，取 $\tau_{m+1}\in(0,\tau_{m})$ 使得对 $\tau_{m+1}<t\leqslant\tau_{m}$ 有 $\xi_{m}(t)<\xi(t)$，且 $\xi_{m}(\tau_{m+1})=\xi(\tau_{m+1})$。显然，该算法经有限步后终止。接着在 $[0,\bar{t}]$ 上构造一个左连续、分段 Lipschitz 的函数 $\xi_{\varepsilon}(\cdot)$：当 $k\geqslant2$ 时它在 $\tau_{1},\cdots,\tau_{k-1}$ 处有跳跃间断；对 $\tau_{m+1}<t\leqslant\tau_{m}$，$m=0,1,\cdots,k-1$，令 $\xi_{\varepsilon}(t)=\xi_{m}(t)$，并令 $\xi_{\varepsilon}(0)=\xi_{k-1}(0)$。于是
\begin{align}
\xi_{\varepsilon}(\bar{t})-\xi_{\varepsilon}(0)=(k-1)\varepsilon+\sum_{m=0}^{k-1}\int_{\tau_{m+1}}^{\tau_{m}}\dot{\xi}_{m}(t)\,\mathrm{d}t
\geqslant\int_{0}^{\bar{t}}\lambda_{i}(U(\xi_{\varepsilon}(t)+,t))\,\mathrm{d}t.
\label{align:ch10-10-3-1}
\end{align}

当 $\varepsilon\downarrow0$ 时，对任意 $t\in[0,\bar{t}]$，$\xi_{\varepsilon}(t)\to\xi(t)$（从左侧）。为看清这一点，固定 $t\in[0,\bar{t}]$。则 $(\xi_{\varepsilon}(t),t)$ 位于由某点 $(\xi(\tau_{\varepsilon})-\varepsilon,\tau_{\varepsilon})$ 发出的极小后向特征线上，其中 $\tau_{\varepsilon}\geqslant t$。取序列 $\{\varepsilon_{k}\}$ 使得 $\varepsilon_{k}\to0$ 且 $\tau_{\varepsilon_{k}}\to\tau$，$k\to\infty$。于是 $(\xi(\tau_{\varepsilon_{k}})-\varepsilon_{k},\tau_{\varepsilon_{k}})\to(\xi(\tau),\tau)$。再由形如\eqref{align:ch10-10-2-1}的偶然方程的标准理论可知，由 $(\xi(\tau_{\varepsilon_{k}})-\varepsilon_{k},\tau_{\varepsilon_{k}})$ 发出的一列极小后向特征线必在 $[0,t]$ 上一致收敛于 $\xi(\cdot)$。特别地，$\xi_{\varepsilon_{k}}(t)\to\xi(t)$，$k\to\infty$。因此在\eqref{align:ch10-10-3-1}中令 $\varepsilon\downarrow0$，得到
\begin{align}
\xi(\bar{t})-\xi(0)\geqslant\int_{0}^{\bar{t}}\lambda_{i}(U(\xi(t)-,t))\,\mathrm{d}t.
\label{align:ch10-10-3-2}
\end{align}

另一方面，在 $[0,\bar{t}]$ 上几乎处处有 $\dot{\xi}(t)\leqslant\lambda_{i}(U(\xi(t)-,t))$；从而在 $[0,\bar{t}]$ 上几乎处处有 $\dot{\xi}(t)=\lambda_{i}(U(\xi(t)-,t))$，亦即 $\xi(\cdot)$ 是一条左 $i$-接触。

同理可证，由 $(\bar{x},\bar{t})$ 发出的极大后向 $i$-特征线是一条右 $i$-接触。证明完毕。
\end{proof}

由\cref{section:ch10-广义特征线}末尾的注记可知，\cref{theorem:ch10-极后向特征线为接触}有如下推论：

\begin{theorem}\label{theorem:ch10-小振幅下极后向特征线无激波}
假设方程组\eqref{align:ch10-10-1-1}的 $i$-特征族真正非线性，且 $U$ 的振幅充分小。则由上半平面任一点 $(\bar{x},\bar{t})$ 发出的极小与极大后向 $i$-特征线均无激波。
\end{theorem}

上述定理的意义将在后续章节中展现。

为后文使用方便，我们在此引入一类由无穷远发出的特殊后向特征线：

\begin{definition}[$i$-分界线与 $i$-分界]\label{definition:ch10-i分界线}
与解 $U$ 相联系的一条\emph{极小（相应地极大）$i$-分界线}，是指一个 Lipschitz 函数 $\xi:[0,\bar{t}]\to(-\infty,\infty)$，使得在紧致时间区间上一致地有 $\xi(t)=\lim_{m\to\infty}\xi_{m}(t)$，其中 $\xi_{m}(\cdot)$ 是由点 $(x_{m},t_{m})$ 发出的极小（相应地极大）后向 $i$-特征线，$t_{m}\to\bar{t}$，$m\to\infty$。特别地，当 $\bar{t}=\infty$ 时，称该 $i$-分界线 $\xi(\cdot)$ 为一条极小（相应地极大）\emph{$i$-分界}。
\end{definition}

注意，任意两条极小（相应地极大）$i$-特征线的图像可能相互相遇，但不能相交。因此，一条极小（相应地极大）后向 $i$-特征线的图像不能与任何极小（相应地极大）$i$-分界线的图像相交。同理，任意两条极小（相应地极大）$i$-分界线的图像也不能相交。特别地，任何极小（相应地极大）$i$-分界都将上半平面分成两部分，使得没有任何前向 $i$-特征线能从左向右（相应地从右向左）穿越它。

极小或极大 $i$-分界线必为广义 $i$-特征线；由\cref{theorem:ch10-极后向特征线为接触}，它们是左或右 $i$-接触。特别地，当 $i$-特征族真正非线性且 $U$ 的振幅充分小时，由\cref{theorem:ch10-小振幅下极后向特征线无激波}可知，极小或极大 $i$-分界线是无激波的。

下列命题给出了分界存在的条件：

\begin{lemma}\label{lemma:ch10-分界存在性}
设存在与 $U$ 相联系、在 $t=0$ 时自 $x$ 轴发出的 $i$-特征线 $\varphi$ 与 $\psi$，满足如下性质：$\varphi(t)<\psi(t)$，$t\in[0,\infty)$，并且 $\psi$ 是极大的（相应地 $\varphi$ 是极小的）。则存在一条极小（相应地极大）$i$-分界 $\xi$，使得
\begin{align*}
\varphi(t)\leqslant\xi(t)\leqslant\psi(t),\qquad t\in[0,\infty).
\end{align*}
\end{lemma}

\begin{proof}
设 $\varphi$ 是极大的，令 $\xi_{m}$ 表示由点 $(x_{m},t_{m})$ 发出的极小后向 $i$-特征线，其中 $\varphi(t_{m})<x_{m}<\psi(t_{m})$，$t_{m}\to\infty$，$m\to\infty$。由于 $\varphi$ 是极大的，$\xi_{m}$ 的图像不能与 $\varphi$ 的图像相交；又由于 $\xi_{m}$ 是极小的，即使它可能与 $\psi$ 的图像相遇，也不能越过它。因此
\begin{align*}
\varphi(t)<\xi_{m}(t)\leqslant\psi(t),\qquad t\in[0,t_{m}].
\end{align*}
由 Ascoli 定理，$\{\xi_{m}\}$ 含子列收敛于一条具有所述性质的极小 $i$-分界 $\xi$。$\psi$ 为极小的情形类似处理。证明完毕。
\end{proof}

\begin{corollary}\label{corollary:ch10-周期解的分界}
当 $U$ 关于空间周期，即 $U(x+L,t)=U(x,t)$ 时，从 $x$ 轴上任一长度为一个周期的区间中，至少发出一条极小 $i$-分界与一条极大 $i$-分界。
\end{corollary}

事实上，为构造极小 $i$-分界，只需对\cref{lemma:ch10-分界存在性}作应用：取 $\varphi$ 为由 $(0,0)$ 发出的极大 $i$-特征线，并令 $\psi(t)=\varphi(t)+L$，$t\in[0,\infty)$。同理，为构造极大 $i$-分界，对\cref{lemma:ch10-分界存在性}作应用：取 $\psi$ 为由 $(L,0)$ 发出的极小 $i$-特征线，再令 $\varphi(t)=\psi(t)-L$，$t\in[0,\infty)$。当然，极小与极大 $i$-分界可能重合。

\begin{remark}[本章结语]\label{remark:ch10-局部性注记}
广义特征线是在与定义于整个上半平面上的\eqref{align:ch10-10-1-1}的 $BV$ 解相联系时引入的。然而显然，这一概念及其许多性质本质上是局部的，因而同样适用于定义在 $\mathbb{R}^{2}$ 任意开子集上的 $BV$ 解。
\end{remark}

\end{document}